% Opening of Chapter 1 (Introduction)

\textit{Yersinia pestis} is the causative agent of bubonic plague. Perhaps more
than any other human pathogen it has struck fear into the hearts and minds of
peoples all over the world, and, consciously or not, it lingers as a shadow in
the western psyche.

This thesis is about the mechanisms the bacterium employs inside a mammalian
host, and about one of them in particular. \yp{} does not resist the immune
system from the outset. It permits itself to be swallowed by the host's own
phagocytes, replicates inside them for a while, and only then reverses course
and refuses uptake altogether. That reversal is among the sharpest differences
between \yp{} and the comparatively mild organism it descends from, and it is
the sort of thing that invites a mathematical answer, because it is a choice
between two strategies whose returns can in principle be written down and
compared.

The picture drawn here is a stochastic one, and it is drawn at the scale where
the choice is actually made: a few bacteria inside a single cell. What such a
population does is settled by chance rather than by its average behaviour, so
the object needed throughout is a distribution and not a trajectory. Most of
the mathematics in this thesis is the machinery that requirement turned out to
need.
